\scanpage{18}
\uline{Théorème 14} Soient $f:X\longrightarrow S$ propre de prés. finie, $F,G$ Modules de prés. finie sur $X$, avec $G$ $S$-plat. Alors, \deleted{\unsure{pour un}} il existe un \deleted{\unsure{Mo}}Module $P$ sur $S$, de présentation finie, et un isom de foncteurs en le $\mathcal O_S$-Module quasi-cohérent $M$
% slanted
\margin{\uline{NB} Il y a \unsure{un même|même un|un nouveau|le même} th. de représentabilité en \unsure{suppt} pas $F$ de prés. fin., mais seulement quasi-cohérent ; \unsure{mais alors|alors} $P$ est seulement qu. coh., et on se ramène au cas $F$ de prés. finie par passage à la limite.}
\[
 \underline{\operatorname{Hom}}(P,M)\simeq
 f_*\underline{\operatorname{Hom}}(F,G\otimes_{\mathcal O_S} M)
\]
(et par suite, en prenant les $\Gamma$, l'isom fonctoriel en $M$
\[
 \operatorname{Hom}(P,M)\simeq\operatorname{Hom}(F,G\otimes_{\mathcal O_S}M),
\]
qui montre que $P$ est déterminé à isom unique près comme représentant le foncteur en $M$ du deuxième mb). La formation de $P$ commute au changement de base quelconques $S'\longrightarrow S$.

On voit facilement qu'on peut supposer $S=\operatorname{Spec}A$ affine, et alors c'est une question de représentabilité d'un foncteur $\operatorname{Mod} A)\longrightarrow(\mathrm{Ens})$. On est ramené comme d'habitude au cas d'un $A$ de type fini sur $\mathbb Z$, donc quotient d'un régulier. On peut donc faire appel au th. 12.

Les conditions (i) à (iv) sont standard. \deleted{\unsure{Pour la}} \inserted{En vertu} de \unsure{remarque} 13, il reste à vérifier la condition (vi).

\deleted{Condition (v), \unsure{on|il} \unsure{est ramené à} prouver que si $A$ est intègre, $\exists$ \deleted{\unsure{ouvert}} ouvert non vide $U$ de $S=\operatorname{Spec}A$ tel que sur $U$, la fonction $s\longrightarrow\dim\operatorname{Hom}(F_s,G_s)$ \deleted{\unsure{$g_s$}} soit constante. Or on a $\operatorname{Hom}(F_s,G_s)=\Gamma(X_s,\underline{\operatorname{Hom}}(F_s,G_s))$, d'autre part pour presque tout}
On peut supposer $B=A$.
% archive sheet number: 17
